% Section 4: Across topics. Written last, from the other sections; each row points back to where the idea is taught.

\cstopic{4.1}{Same law, three settings}{1.1, 1.2, 3.1}
\takeaway{This topic lines up the three places the course uses the decay law: counting nuclei, measuring activity and simulating decay. Questions that mix units usually hinge on noticing that all three halve at the same rate.}

\begin{tablebox}{Nuclei, activity, simulation}{@{}>{\raggedright\arraybackslash}p{0.2\linewidth}>{\raggedright\arraybackslash}p{0.27\linewidth}>{\raggedright\arraybackslash}p{0.23\linewidth}X@{}}
\toprule & \textbf{Quantity} & \textbf{Law} & \textbf{Halves every}\\\midrule
Nuclei (1.1) & \(N(t)\) & \(N_0 e^{-\lambda t}\) & \(t_{1/2}=\frac{\ln 2}{\lambda}\)\\
Activity (1.2) & \(A(t)=\lambda N\) & \(A_0 e^{-\lambda t}\) & the same \(t_{1/2}\)\\
Simulation (3.1) & expected \(N_k\) & \(N_0(1-p)^k\) & about \(\frac{\ln 2}{p}\) steps\\
\bottomrule
\end{tablebox}

\begin{tfbox}
\tf{T}{If a detector's count rate halves in 3 hours, the number of nuclei also halves in 3 hours.}{\(A=\lambda N\).}
\tf{F}{A larger sample takes longer to lose half its nuclei.}{The half-life does not depend on \(N_0\).}
\end{tfbox}
